\begin{frame}{Caso 1. Razón corriente y prueba ácida}
\smallcase
\begin{columns}[T,totalwidth=\textwidth]
\column{0.58\textwidth}
\casehead{Enunciado.} Comercial Andina S.A.C., distribuidora ubicada en Lima, desea evaluar su capacidad para atender obligaciones de corto plazo.

\casehead{Datos.}
Activo corriente = \soles{480,000.00}; Inventarios = \soles{180,000.00}; Pasivo corriente = \soles{300,000.00}.

\casehead{Operación.}
\[
\text{Razón corriente}=\frac{\soles{480,000.00}}{\soles{300,000.00}}=1.60
\]
\[
\text{Prueba ácida}=\frac{\soles{480,000.00}-\soles{180,000.00}}{\soles{300,000.00}}=1.00
\]
\casehead{Respuesta.} Razón corriente = 1.60; prueba ácida = 1.00.

\casehead{Interpretación.} Por cada \soles{1.00} de deuda corriente dispone de \soles{1.60} en activos corrientes. Al excluir inventarios, la cobertura disminuye a \soles{1.00}.

\column{0.37\textwidth}
\centering
\begin{tikzpicture}[node distance=0.48cm, every node/.style={font=\small}]
\node[draw, rounded corners=3pt, fill=liquidez!10, align=center, minimum width=3.2cm] (a) {Con inventarios\\\textbf{1.60}};
\node[draw, rounded corners=3pt, fill=liquidez!5, align=center, minimum width=3.2cm, below=of a] (b) {Sin inventarios\\\textbf{1.00}};
\draw[-{Stealth}, thick] (a) -- (b);
\end{tikzpicture}
\end{columns}
\end{frame}
